\section{Literature Review}

This section evaluates the theoretical foundations and state-of-the-art methodologies underlying the proposed PyTrade architecture. The system separates temporal signal extraction from cross-sectional portfolio allocation. This review systematically establishes the necessity of each architectural component, transitioning from foundational machine learning principles to specialized quantitative finance applications.

\subsection{Sequential Decision-Making and Continuous Control in Finance}
Reinforcement learning models sequential decision-making under uncertainty as a Markov Decision Process (MDP), defined by states, actions, state-transition probabilities, and scalar rewards. Goodfellow et al. \cite{Goodfellow-et-al-2016} document how deep neural networks act as high-capacity function approximators, enabling policy networks to learn complex mapping functions directly from high-dimensional input spaces. In continuous action domains, traditional policy gradient methods suffer from high variance and sample inefficiency. To resolve this, Schulman et al. \cite{JohnSchulmanFilipWolskiPrafullaDhariwalAlecRadfordOlegKlimov.2017} introduced Proximal Policy Optimization (PPO), which clips the surrogate objective function during gradient updates. This clipping mechanism bounds policy ratio updates, preventing destructive step changes and stabilizing convergence in continuous state-action environments.

In quantitative finance, portfolio management represents a continuous control problem where trading decisions alter available capital, market exposures, and transaction overhead over time. Early financial machine learning frameworks often simplified execution by discretizing decisions into binary buy, sell, or hold events. To formalize discrete execution boundaries, López de Prado \cite{LopezdePrado.2018} introduced the triple-barrier method, which labels price paths based on upper take-profit limits, lower stop-loss thresholds, and maximum holding durations. However, discrete actions cannot dynamically size capital allocations across multi-asset portfolios. Addressing this limitation, Wu et al. \cite{Junfeng.2024} demonstrated the effectiveness of PPO in continuous action spaces for dynamic portfolio management, utilizing shared representation layers to extract general market features before splitting into specialized actor and critic heads. To standardize benchmark evaluations in this domain, Liu et al. \cite{liu2022finrldeepreinforcementlearning} introduced the FinRL library, establishing unified environments and execution pipelines for financial DRL. Further expanding continuous control objectives, Charkhestani and Esfahanipour \cite{Charkestani.2026} incorporated behavioral finance principles into continuous policy optimization, demonstrating that loss aversion and overconfidence constraints can be embedded directly into agent reward structures to regulate risk-taking behavior.

\subsection{Sequential Representation and Temporal Memory in Non-Stationary Markets}
Financial markets are non-stationary and partially observable systems where observable price feeds fail to capture complete macroeconomic state variables or private order flows. When state representation is incomplete, the standard Markov assumption breaks down. To solve Partially Observable MDPs (POMDPs), Hausknecht and Stone \cite{DBLP:journals/corr/HausknechtS15} established Deep Recurrent Q-Learning, integrating recurrent memory layers to maintain internal state representations over sequence steps. Hochreiter and Schmidhuber \cite{Hochreiter1997} established the Long Short-Term Memory (LSTM) architecture, introducing specialized gating mechanisms (input, forget, and output gates) that regulate information flow to eliminate vanishing gradients and capture long-range temporal dependencies.

Applying sequential models to financial data introduces a fundamental trade-off between statistical stationarity and temporal memory. Statistical stationarity, where the mean and variance of a time series remain constant over time, is essential for stable neural network training. Standard integer differencing (e.g., computing log returns) achieves stationarity but erases historical price levels, removing persistent macro trends. To resolve this conflict, López de Prado \cite{LopezdePrado.2018} advocated for fractional differentiation, an approach that applies non-integer differencing operators to satisfy statistical stationarity while preserving maximum historical memory. In parallel, advanced sequence architectures have expanded time-series forecasting capability. Lim et al. \cite{lim2020tft} introduced Temporal Fusion Transformers (TFT), combining recurrent layers with temporal self-attention for interpretable multi-horizon forecasting. Similarly, Shen et al. \cite{LefeiShen.2025} conducted a deep evaluation of Transformer architectures, establishing their capacity to model long-term sequence dependencies in complex time-series data. Despite these forecasting advances, recurrent state representations remain the primary choice for online continuous control policies due to their low-latency state updating.

\subsection{Self-Attention and Cross-Sectional Asset Interaction}
Attention mechanisms calculate dynamic weighting factors across input features based on relative contextual relevance. Vaswani et al. \cite{vaswani2023attentionneed} introduced the Transformer architecture, replacing recurrent sequence bottlenecks with multi-head self-attention. Self-attention maps input vectors into Query, Key, and Value matrices, allowing every element in a sequence or set to compute pairwise interactions with all other elements simultaneously. Hu et al. \cite{Hu.2024} surveyed the trajectory of Transformer-based reinforcement learning, documenting how attention models handle multi-modal context, trajectory history, and spatial representation across diverse decision tasks. Expanding continuous control capabilities, Kachaev et al. \cite{NikitaKachaev.2025} demonstrated that Transformer encoders enhance representation learning in online continuous control settings, enabling agents to parse non-stationary state representations efficiently.

In multi-asset portfolio management, an asset's risk and return dynamics cannot be evaluated in isolation; potential allocations depend on cross-asset correlations, asset dispersion, and portfolio-wide risk limits. Traditional quantitative methods rely on empirical covariance matrices, which degrade rapidly during market regime shifts. Attention mechanisms solve this by learning non-linear cross-asset dependencies directly from data. Xue and Ye \cite{Xue.2026} demonstrated attention-enhanced reinforcement learning for dynamic portfolio optimization, validating that self-attention layers successfully model transaction cost frictions and multi-asset risk metrics. Furthermore, Coulter \cite{Coulter2025} proposed neural portfolio allocators that process distinct asset representations via multi-head attention to evaluate cross-sectional allocation opportunities dynamically.

\subsection{Modular Decomposition and Hierarchical Policy Architectures}
Monolithic deep reinforcement learning agents often experience policy collapse when forced to solve temporal feature extraction, cross-sectional covariance modeling, and portfolio constraint satisfaction simultaneously within a single neural network \cite{Hu.2024}, \cite{JunKevin.2025}. Hierarchical Reinforcement Learning (HRL) mitigates this optimization complexity by decomposing tasks into specialized, modular sub-policies. In financial domains, Kevin and Yugopuspito \cite{JunKevin.2025} implemented a hybrid LSTM-PPO framework for dynamic portfolio optimization, demonstrating that decoupling recurrent temporal features from policy optimization improves learning stability. Addressing non-stationary market conditions, Raj \cite{GabrielNixonRaj.2025} formulated adaptive, regime-aware RL models, emphasizing that specializations across distinct volatility regimes enhance agent robustness. Complementing regime-aware modeling, Bennett \cite{Bennett.2014} established that financial asset classes exhibit distinct term structures, volatility skew, and cross-correlation patterns that require localized risk management strategies.

The literature demonstrates a possible trajectory toward modularity, continuous control, and attention-based cross-sectional modeling. However, combining single-asset temporal feature extraction with multi-asset attention mechanisms within a unified hierarchical DRL framework remains underexplored. The proposed PyTrade architecture addresses this gap by decoupling single-asset temporal signals from portfolio-level allocation, establishing the theoretical foundation for the system design detailed in Section \ref{sec:system-design}.